\section{A procedure has preconditions and effects}\label{sec:C15:procedure}
A skill is reusable procedural knowledge that changes how an agent approaches a class of tasks. It can be a natural-language instruction, executable code, a workflow, or a package combining these forms. Its defining feature is a claim about applicability and action, rather than a claim about a fact alone. ``The repository uses a particular test framework'' is semantic knowledge; ``run this fixture before changing the parser'' is procedural guidance. A deployed skill can contain both, but the distinction matters when its environment changes.

Represent a candidate skill as $s=(P,\rho,Q,E,V)$, where $P$ is an applicability predicate, $\rho$ the procedure or policy fragment, $Q$ an intended postcondition, $E$ the evidence supporting it, and $V$ compatibility information. This is a proposed contract, not a requirement that every skill system exposes these fields. Natural-language skills often leave $P$ and $Q$ implicit. That makes them flexible but creates a risk that the policy applies a procedure beyond the conditions under which its evidence was collected.

Procedural memory is external to the model when it is stored and retrieved as an artifact. Activating the skill changes the model's conditional input; it does not necessarily change model parameters. A skill bank can improve performance with a frozen backbone, but that result should not be described as weight learning. Conversely, a procedure can later be distilled into parameters, changing its update, deletion, and verification boundaries. Chapter~\ref{ch:C30} treats the difference between these improvement targets.

The Agent Skills specification describes an artifact interface for instructions and resources~\cite{agentskills2026spec}. It is an interoperability and packaging reference, not empirical evidence that an agent chooses or executes skills correctly. A tool schema, a skill package, and a model parameter update are therefore distinct objects. A tool exposes an executable capability; a skill can tell the agent when or how to use capabilities; learning changes the policy that interprets both. Treating them as interchangeable makes attribution and authority ambiguous.

The institutional-knowledge perspective emphasizes that reusable procedures can encode local practices and constraints~\cite{corpus260314805}. Its case-study evidence should not be mistaken for a controlled estimate of productivity. The more general engineering implication is to attach scope and provenance to a skill rather than treating every successful local workaround as universally reusable knowledge. A workaround for one model's failure may become harmful after the model or tool interface is upgraded.

